%O frontmmatter são as chamadas páginas iniciais, terá de atualizar as respetivas secções.

%-----------------------------------------%
%    CONFIGURAÇÃO INICIAL                 %
%-----------------------------------------%

\input{frontmatter/glossary}

\frontmatter % Use roman page numbering style (i, ii, iii, iv...) for the pre-content pages

\pagestyle{plain} % Default to the plain heading style until the thesis style is called for the body content

%----------------------------------------------------------------------------------------
%	CAPA
%----------------------------------------------------------------------------------------
\maketitlepage


%----------------------------------------------------------------------------------------
%	CONTRA CAPA
%----------------------------------------------------------------------------------------
\makeconttitlepage


%-----------------------------------------%
%      EPÍGRAFE     (opcional)                     
%-----------------------------------------%
\begin{epigraph}
\null \vfill
\begin{flushright}

“The reasonable man adapts himself to the world: the unreasonable one persists in trying to adapt the world to himself. Therefore all progress depends on the unreasonable man.”

George Bernard Shaw

\end{flushright}
\vfill \null%
\end{epigraph}%



%----------------------------------------------------------------------------------------
%	DEDICATÓRIA  (opcional)
%----------------------------------------------------------------------------------------
%\dedicatory{For/Dedicated to/To my\ldots}
%\begin{dedicatory}
%\null \vfill
%\begin{flushright}
 

%\end{flushright}
%\vfill \null%
%\end{dedicatory}


%----------------------------------------------------------------------------------------
%	AGRADECIMENTOS (opcional)
%----------------------------------------------------------------------------------------

%\begin{acknowledgements}

%\end{acknowledgements}




%----------------------------------------------------------------------------------------
%	RESUMO
%----------------------------------------------------------------------------------------

%\begin{abstract}

%\noindent \textbf{\ Subtítulo caso queira!}

%Segue-se, com caráter obrigatório, um resumo em língua portuguesa e em língua inglesa (abstract), cada um deles com um máximo de 300 %palavras.

%Após cada um dos resumos devem ser indicadas cinco palavras-chave – %português e inglês – para indexação futura.

% As palavras chave terão de ser definidas no ficheiro main.tex depois da linha de código keywords
%\end{abstract}


%----------------------------------------------------------------------------------------
%	ABSTRACT 
%----------------------------------------------------------------------------------------
%\begin{abstractotherlanguage}
% here you put the abstract in the "other language": English, if the work is written in Portuguese; Portuguese, if the work is written in English.


%\noindent \textbf{\ Subtitle if you want}
%\noindent 

%Trabalhos escritos em língua Inglesa devem incluir um resumo alargado com cerca de 1000 palavras, ou duas páginas.

%Se o trabalho estivesse escrito em Português, este resumo seria em língua Inglesa, com cerca de 200 palavras, ou uma página.

%Para alterar a língua basta ir às configurações do documento no ficheiro \file{main.tex} e alterar para a língua desejada ('english' ou 'portuguese')\footnote{Alterar a língua requer apagar alguns ficheiros temporários; O target \keyword{clean} do \keyword{Makefile} incluído pode ser utilizado para este propósito.}. Isto fará com que os cabeçalhos incluídos no template sejam traduzidos para a respetiva língua.

% As palavras chave terão de ser definidas no ficheiro main.tex depois da linha de código conkeywords

%\end{abstractotherlanguage}
\checktoopen



%----------------------------------------------------------------------------------------
%	ÍNDICE DE CONTEÚDO / FIGURAS / TABELAS
%----------------------------------------------------------------------------------------

\tableofcontents % Imprime o índice principal
\pdfbookmark[0]{\contentsname}{toc}% Adiciona o índice aos bookmarks do pdf
\checktoopen

%----------------------------------------------------------------------------------------
%	ABREVIATURAS
%----------------------------------------------------------------------------------------
\chapter*{\abbrevname}
\addcontentsline{toc}{chapter}{\abbrevname} % Adiciona a lista ao índice principal
%List of Abreviations
%
\begin{acronym}[GPOS] % Coloque aqui a sigla mais comprida para alinhar a tabela
    \acro{BLDC}{\emph{Brushless DC}}
    \acro{WSF}{What (it) Stands For}
    \acro{RTOS}{Real-Time Operating System}
    \acro{GPOS}{General Purpose Operating System}
\end{acronym}
\checktoopen

\listoffigures % Imprime a lista de figuras
\checktoopen

\listoftables % Imprime a lista de tabelas
\checktoopen

\listofmyequations  %Imprime a lista de equações
\checktoopen

%\listofalgorithms % Prints the list of algorithms
%\addchaptertocentry{\listalgorithmname} %Uncomment para mostrar no índice a lista de algoritmos

%\lstlistoflistings % Imprime a lista de listagens (código-fonte da linguagem de programação)

%\addchaptertocentry{\lstlistlistingname} %Uncomment para mostrar a lista de de código no índice





%----------------------------------------------------------------------------------------
%	SÍMBOLOS
%----------------------------------------------------------------------------------------

%\begin{symbols}{lll} % Inclui uma lista de símbolos (uma tabela de três colunas)

%\input{frontmatter/listsymbols}

%\end{symbols}



%----------------------------------------------------------------------------------------
%	ACRÓNIMOS
%----------------------------------------------------------------------------------------

%Use GLS
%\glsresetall
%\printglossary[title=\listacronymname,type=\acronymtype,style=long]


%----------------------------------------------------------------------------------------
%	ACABOU - BOM TRABALHO
%----------------------------------------------------------------------------------------

\mainmatter % Começar numeração da página com numéros árabes (1,2,3 ...)

